\clearpage
\appendices
\section{Notation}
\label{app:notation}

\begin{table}[!t]
\caption{Main symbols used in the forecasting formulation.}
\label{tab:notation}
\centering
\small
\begin{tabularx}{\columnwidth}{@{}l>{\raggedright\arraybackslash}X@{}}
\toprule
Symbol & Meaning \\
\midrule
$G(t)$ & Measured irradiance channel in the historical dataset \\
$\widehat{G}_{c}(t)$ & Clipped Stage-1 estimated irradiance proxy used by the proposed forecast path \\
$\Gcs(t)$ & Clear-sky irradiance reference \\
$\kt(t)$ & Clearness-index-like transmissivity signal \\
$\widehat{\kt}_{c}(t)$ & Clipped Stage-1 transmissivity estimate \\
$\Ppv(t)$ & PV plant AC power output \\
$\Prated$ & Rated plant capacity \\
$P_{\mathrm{cs}}(t)$ & Clear-sky power feature derived from $\Gcs(t)$ \\
$\texttt{theoretical\_power}$ & Logged historical reference channel, not the proposed path's target-time envelope input \\
$A_{\mathrm{eff}}$ & Effective irradiance-to-power coefficient, 181.591 kW per kW/m$^2$ in the active metadata \\
$\widehat{P}_{\mathrm{theo},c}(t)$ & Nonnegative estimated theoretical-power envelope derived from $\widehat{G}_{c}(t)$ \\
$\alpha(t)$ & Solar elevation angle \\
$\omega(t)$ & Solar hour angle \\
$\delta_n$ & Solar declination for day $n$ \\
$\hat{y}$ & Forecast value \\
\bottomrule
\end{tabularx}
\end{table}

Table~\ref{tab:notation} separates measured irradiance, estimated irradiance, clear-sky irradiance, raw logged theoretical power, and plant-power quantities because these variables have different roles in the evaluation. The proposed forecast path uses $\widehat{G}_{c}(t)$ and $\widehat{P}_{\mathrm{theo},c}(t)$ derived from Stage 1, while the sensor-assisted diagnostic path uses measured irradiance-derived features at the target timestamp. This separation matters when comparing the proposed path with persistence and sensor-assisted baselines.

\section{Feature-Engineering Trace}
\label{app:feature_trace}

The feature pipeline first converts each timestamp into hour of day, day of year, and month. These temporal variables are represented with sine and cosine harmonics, and the day-of-year signal also includes a second harmonic. This encoding represents daily and seasonal cycles smoothly for the tabular tree model.

Solar features are computed from latitude, day of year, and civil hour using the radian equations in Section~\ref{sec:methodology}. The implementation uses Indian Standard Time directly rather than local apparent solar time. The resulting solar-elevation angle, after inverse-sine clamping, is used to compute a clear-sky irradiance reference, the air-mass-like feature in Eq.~\eqref{eq:airmass}, a clear-sky power feature, and a 1-h smoothed clear-sky irradiance feature. These variables provide a reproducible physical envelope before cloud and atmospheric attenuation are modeled.

Meteorological variables are expanded into anomaly, humidity, dew-point, and rolling-window features. Temperature and humidity anomalies are computed relative to monthly means. Dew point follows the simplified expression in Eq.~\eqref{eq:dewpoint}. Precipitation is accumulated over 1-h, 3-h, and 24-h windows using Eq.~\eqref{eq:rolling_sum}, while humidity is summarized over a 3-h rolling window using Eq.~\eqref{eq:rolling}. These derived variables describe slowly changing atmospheric conditions without introducing a recurrent model.

Cloud variables enter through both raw layer channels and derived total-cloud features. Total, low, mid, and high cloud-cover channels are passed to the Stage-1 model. Total cloud cover also contributes a 1-h lag, a 3-h rolling mean, and a 1-h change feature. This design allows the model to distinguish a stable overcast sky from a rapidly changing cloud field.

The sensor-assisted diagnostic model and the proposed forecast path differ in their resource inputs. The diagnostic path uses measured irradiance-derived rolling features and the raw \texttt{theoretical\_power} reference channel at the target timestamp. The proposed path instead uses clipped estimated irradiance features produced by Stage 1 and the derived envelope in Eq.~\eqref{eq:theo_est}. The proposed path therefore evaluates how well meteorological and cloud variables can reconstruct a solar-resource signal, while the sensor-assisted path estimates the plant-power mapping when the resource signal is already known.

\section{Evaluation Masks and Metric Traceability}
\label{app:metric_trace}

All headline metrics use the chronological held-out test interval and the post-processing rule in Eq.~\eqref{eq:postprocess}. The headline daylight mask is defined by $\Gcs>10$ W/m$^2$, which yields the proposed-path test set reported in Table~\ref{tab:headline_metrics}. Table~\ref{tab:evaluation_masks} gives the full mask trace, including the $\alpha(t)>5^\circ$ Stage-1 training mask and the feature-availability filters. This daylight mask avoids allowing night intervals to dominate the normalized error metrics.

Baseline skill is reported only on common masks. The previous-day persistence comparison removes rows where the 24-h lag is unavailable, while the 5-min last-value comparison requires $y(t-5~\mathrm{min})$. These common-mask definitions explain why the row counts in Table~\ref{tab:baseline_summary} differ slightly across baseline comparisons. Mixing masks, including all night rows, or comparing baselines on different samples would change the reported skill values.

The ablation comparisons use the same chronological split and the same daylight metric convention as the proposed path. The single-stage model removes the learned transmissivity intermediate, while the total-cloud-only cascade removes the low-, mid-, and high-cloud raw channels. This matched evaluation isolates the effect of model structure and cloud-layer availability rather than differences in test period or row filtering.

\section{Forecast Horizons}
\label{app:horizon_note}

The reported experiment evaluates predictions on a 5-min target grid, but the operational meaning of that grid depends on feature availability. In this manuscript, the held-out test uses target-time meteorological and cloud columns from the same retained dataset, so the reported RMSE is an offline target-weather-input result. If the same target-time fields are supplied by a forecast service or local weather pipeline before issue time, the model becomes a weather-input forecast for those target timestamps. If measured irradiance is supplied at the target timestamp, the task becomes a sensor-assisted plant mapping. If target-adjacent plant power is supplied, a last-value or autoregressive model becomes a live nowcast.

For this reason, the proposed weather-input path should not be interpreted as a universal replacement for a one-step live nowcast. The 5-min last-value persistence baseline has access to plant output immediately before the target timestamp, which gives it stronger short-horizon information than the proposed weather-input path. The positive skill against previous-day same-time persistence is the more relevant comparison for schedule-oriented or weather-input operation under the present protocol.

Future horizon-specific experiments should report separate results for one-step-ahead, 30-min-ahead, 1-h-ahead, and day-ahead settings. Very short horizons should be compared with last-value, ramp-aware persistence, and autoregressive baselines that use recent plant power. Longer horizons should be compared with previous-day, clear-sky-index persistence, climatology, and weather-forecast-driven baselines. This separation keeps the evaluation aligned with the information actually available at each forecast horizon.

\section{CSV Inference}
\label{app:forecast_csv_demo}

The trained pipeline also includes a forecast-CSV inference path. In this mode, a weather-input file is normalized, cloud-layer fields are used directly or estimated from the month-hour climatology, hourly inputs are interpolated to the 5-min grid, and the same Stage-1 and Stage-2 models are applied with the post-processing rule in Eq.~\eqref{eq:postprocess}. This trial demonstrates the inference workflow; it is not used as a validation metric because measured PV output for the same forecast window is not available.

\begin{figure}[!t]
\centering
\includegraphics[width=\columnwidth]{figures/results/forecast_csv_inference_power.png}
\caption{Forecast-CSV inference demonstration using external weather inputs. The curve shows the 5-min power profile generated from the weather-input CSV; no measured PV power is available for this window, so the figure is a workflow demonstration rather than an error-metric result.}
\label{fig:forecast_csv_demo}
\end{figure}

\begin{table}[!t]
\caption{Forecast-CSV inference demonstration summary.}
\label{tab:forecast_csv_demo}
\centering
\small
\begin{tabularx}{\columnwidth}{@{}>{\raggedright\arraybackslash}X>{\raggedright\arraybackslash}X@{}}
\toprule
Quantity & Demonstration result \\
\midrule
Input file & Root-active weather-input CSV example \\
Weather-input window & 2026-05-14 00:00 to 2026-05-27 23:00 \\
Generated prediction grid & 4,020 rows at 5-min cadence \\
Calendar dates represented & 14 \\
Total predicted energy & 3,319.0 kWh \\
Mean daily predicted energy & 237.1 kWh/day \\
Mean daily peak power & 34.4 kW \\
Highest-energy day & 2026-05-16, 274.2 kWh \\
Lowest-energy day & 2026-05-14, 191.2 kWh \\
Predicted hours above 30 kW & 42 h \\
\bottomrule
\end{tabularx}
\end{table}


The demonstration window covers 14 calendar dates from 2026-05-14 to 2026-05-27. The model produces 4,020 5-min predictions, a total predicted energy of 3,319.0 kWh, and a mean daily predicted energy of 237.1 kWh/day. These values show that the saved models, feature logic, and envelope clipping can be reused outside the training notebook once the weather-input schema is satisfied. They should not be compared with the held-out test metrics in Table~\ref{tab:headline_metrics}, because the demonstration has no same-window measured plant-power target.
